\documentclass[11pt]{amsart}

\usepackage[T1]{fontenc}
\usepackage{lmodern}
\usepackage{microtype}
\usepackage{amsmath,amssymb}
\usepackage{booktabs}
\usepackage[hidelinks]{hyperref}

\newtheorem{theorem}{Theorem}

\newcommand{\Pten}{P_{10}}
\newcommand{\Ptwelve}{P_{12}}
\newcommand{\Ba}{B^{(a)}}
\newcommand{\Bb}{B^{(b)}}
\newcommand{\uv}[2]{#1\!-\!#2}
\newcommand{\doi}[1]{\href{https://doi.org/#1}{doi:#1}}

\hypersetup{
  pdftitle={The Blanusa Snarks Disprove the P12-Conjecture},
  pdfauthor={Owen Shen},
  pdfsubject={H-colorings of cubic graphs},
  pdfkeywords={cubic graph, H-coloring, Petersen coloring, perfect matching, snark}
}

\title{The Blanu\v{s}a Snarks Disprove the
\texorpdfstring{$P_{12}$}{P12}-Conjecture}
\author{Owen Shen}
\address{Operations Research Center, Massachusetts Institute of Technology,
Cambridge, MA 02139, USA}
\email{owenshen@mit.edu}
\date{August 2026}

\subjclass[2020]{05C15, 05C70}
\keywords{$H$-coloring, Petersen coloring, perfect matching index, cubic graph,
snark}

\begin{document}

\begin{abstract}
Let $P_{12}$ be obtained from the Petersen graph by expanding one vertex to a
triangle. Hakobyan and Mkrtchyan conjectured that every cubic graph $G$ with
perfect-matching index $k(G)\leq4$ admits a $P_{12}$-coloring. We prove that
the two $18$-vertex Blanu\v{s}a snarks are counterexamples: each has
$k(G)=4$, admits a Petersen coloring, and admits no $P_{12}$-coloring. Thus
Conjecture~8 is false, and $P_{12}$-colorability is strictly stronger than
Petersen-colorability even within the class $k(G)=4$. The nonexistence
assertions are certified by exact exhaustive enumeration.
\end{abstract}

\maketitle

\section{The counterexamples}

All graphs are finite and loopless. For a vertex $v$ of a graph $G$, write
$\partial_G(v)$ for the set of edges incident with $v$. If $G$ and $H$ are
cubic graphs, an \emph{$H$-coloring} of $G$ is a map
$f:E(G)\to E(H)$ such that, for every $v\in V(G)$, there is an
$x\in V(H)$ satisfying
\[
  f\bigl(\partial_G(v)\bigr)=\partial_H(x).
\]
We write $H\prec G$ when such a coloring exists. This relation is transitive
\cite{HakobyanMkrtchyan2019}.

Let $k(G)$ be the minimum number of perfect matchings whose union is $E(G)$,
with $k(G)=\infty$ if no such cover exists. Let $\Pten$ denote the Petersen
graph, and let $\Ptwelve$ be obtained from $\Pten$ by replacing one vertex
with a triangle and attaching its three former incident edges to distinct
vertices of that triangle. Hakobyan and Mkrtchyan posed the following as
Conjecture~8 \cite{HakobyanMkrtchyan2019}:
\begin{equation}\label{eq:HM}
  k(G)\leq 4 \quad\Longrightarrow\quad \Ptwelve\prec G
  \qquad\text{for every cubic graph }G.
\end{equation}

Let $\Ba$ and $\Bb$ denote the two Blanu\v{s}a snarks, with the exact labeled
edge sets specified in Appendix~\ref{app:data}; the superscripts are only
labels for the present data.

\begin{theorem}\label{thm:main}
For each $X\in\{\Ba,\Bb\}$,
\[
  k(X)=4,\qquad \Pten\prec X,\qquad\text{and}\qquad
  \Ptwelve\not\prec X.
\]
Consequently, Conjecture~8 of Hakobyan and Mkrtchyan is false.
\end{theorem}

\begin{proof}
For each $X\in\{\Ba,\Bb\}$, the four perfect matchings displayed in
Appendix~\ref{app:data} cover $E(X)$, so $k(X)\leq4$.

The remaining assertions are checked by exact finite enumeration. For either
simple cubic target $H\in\{\Pten,\Ptwelve\}$, a local state at a source vertex
$v$ consists of a vertex $x\in V(H)$ and a bijection
$\partial_X(v)\to\partial_H(x)$. Two states at adjacent source vertices are
compatible when they assign the same target edge to their common source edge.
Thus $H$-colorings are exactly the compatible choices of one local state at
each source vertex. Proper $3$-edge-colorings are enumerated analogously by
using the six bijections from each source star to $\{1,2,3\}$.

A depth-first search with exact domain propagation gives the following numbers
of labeled colorings to fixed labeled targets:
\begin{center}
\begin{tabular}{lrrr}
\toprule
source & proper $3$-edge-colorings & $\Pten$-colorings & $\Ptwelve$-colorings\\
\midrule
$\Ba$ & 0 & 480 & 0\\
$\Bb$ & 0 & 240 & 0\\
\bottomrule
\end{tabular}
\end{center}
Hence $\Pten\prec X$ and $\Ptwelve\not\prec X$ for both graphs.

Finally, no cover by fewer than three perfect matchings is possible in a
cubic graph. A cover by three exists exactly when the graph is properly
$3$-edge-colorable: the color classes give one direction, while three perfect
matchings contain $3|V(G)|/2=|E(G)|$ edge occurrences, so any such cover must
partition $E(G)$. Since neither graph is $3$-edge-colorable, the displayed
four-covers are minimum, and $k(X)=4$.

The accompanying standard-library verifier stores the labeled graph data,
checks the matching covers, constructs $\Pten$ and $\Ptwelve$ from their
definitions, and reproduces the complete enumeration above.
\end{proof}

Hakobyan and Mkrtchyan observed that $\Pten\prec\Ptwelve$
\cite{HakobyanMkrtchyan2019}. By transitivity, every $P_{12}$-colorable cubic
graph is Petersen-colorable. Theorem~\ref{thm:main} proves that the converse
fails even when the perfect-matching index is four.

\section{Context}

Blanu\v{s}a introduced the two graphs in 1946 \cite{Blanusa1946}. Fouquet and
Vanherpe proved that both have perfect-matching index four
\cite{FouquetVanherpe2009}, and H\"agglund and Steffen proved that all
generalized Blanu\v{s}a snarks of both types admit Petersen colorings
\cite{HagglundSteffen2014}. The contribution of this note is the
non-$P_{12}$-colorability of the two classical members, which both refutes
\eqref{eq:HM} and separates $P_{12}$-colorability from Petersen-colorability
inside the class $k(G)=4$.

Putman recently exhibited a $112$-vertex simple bridgeless cubic graph with no
Petersen coloring \cite{Putman2026}. The ancillary material includes an explicit
four-perfect-matching cover of his graph and a standard-library checker. The
cover gives $k(G)\leq4$, while non-Petersen-colorability rules out a proper
$3$-edge-coloring and hence a cover by three perfect matchings. Thus Putman's
graph also has $k(G)=4$ and is another counterexample to Conjecture~8. The
examples in Theorem~\ref{thm:main} have a different and sharper profile: they
have order $18$ and remain Petersen-colorable.

\appendix
\section{Labeled graphs and perfect-matching covers}\label{app:data}

Both vertex sets are $\{0,1,\ldots,17\}$. We write $u-v$ for the edge with
endpoints $u$ and $v$.

\paragraph{Edge sets.}
\begingroup\footnotesize
\begin{align*}
E(\Ba)=\{&
\uv{0}{4},\uv{0}{5},\uv{0}{12},\uv{1}{2},\uv{1}{6},\uv{1}{13},
\uv{2}{7},\uv{2}{10},\uv{3}{4},\\
&\uv{3}{8},\uv{3}{14},\uv{4}{9},\uv{5}{7},\uv{5}{8},\uv{6}{8},
\uv{6}{9},\uv{7}{9},\uv{10}{11},\\
&\uv{10}{15},\uv{11}{12},\uv{11}{16},\uv{12}{17},\uv{13}{15},
\uv{13}{16},\uv{14}{16},\uv{14}{17},\uv{15}{17}\},
\end{align*}
\begin{align*}
E(\Bb)=\{&
\uv{0}{4},\uv{0}{5},\uv{0}{12},\uv{1}{2},\uv{1}{6},\uv{1}{13},
\uv{2}{3},\uv{2}{7},\uv{3}{4},\\
&\uv{3}{10},\uv{4}{9},\uv{5}{7},\uv{5}{8},\uv{6}{8},\uv{6}{9},
\uv{7}{9},\uv{8}{14},\uv{10}{11},\\
&\uv{10}{15},\uv{11}{12},\uv{11}{16},\uv{12}{17},\uv{13}{15},
\uv{13}{16},\uv{14}{16},\uv{14}{17},\uv{15}{17}\}.
\end{align*}
\endgroup

\paragraph{Four-perfect-matching covers.}
For $\Ba$, take
\begingroup\footnotesize
\begin{align*}
M^{(a)}_1={}&\{\uv{0}{4},\uv{1}{2},\uv{3}{8},\uv{5}{7},\uv{6}{9},
\uv{10}{11},\uv{12}{17},\uv{13}{15},\uv{14}{16}\},\\
M^{(a)}_2={}&\{\uv{0}{4},\uv{1}{6},\uv{2}{10},\uv{3}{14},\uv{5}{8},
\uv{7}{9},\uv{11}{12},\uv{13}{16},\uv{15}{17}\},\\
M^{(a)}_3={}&\{\uv{0}{5},\uv{1}{13},\uv{2}{7},\uv{3}{14},\uv{4}{9},
\uv{6}{8},\uv{10}{15},\uv{11}{16},\uv{12}{17}\},\\
M^{(a)}_4={}&\{\uv{0}{12},\uv{1}{6},\uv{2}{10},\uv{3}{4},\uv{5}{8},
\uv{7}{9},\uv{11}{16},\uv{13}{15},\uv{14}{17}\}.
\end{align*}
\endgroup
For $\Bb$, take
\begingroup\footnotesize
\begin{align*}
M^{(b)}_1={}&\{\uv{0}{4},\uv{1}{2},\uv{3}{10},\uv{5}{7},\uv{6}{9},
\uv{8}{14},\uv{11}{12},\uv{13}{16},\uv{15}{17}\},\\
M^{(b)}_2={}&\{\uv{0}{4},\uv{1}{6},\uv{2}{3},\uv{5}{8},\uv{7}{9},
\uv{10}{11},\uv{12}{17},\uv{13}{15},\uv{14}{16}\},\\
M^{(b)}_3={}&\{\uv{0}{5},\uv{1}{13},\uv{2}{7},\uv{3}{4},\uv{6}{9},
\uv{8}{14},\uv{10}{15},\uv{11}{16},\uv{12}{17}\},\\
M^{(b)}_4={}&\{\uv{0}{12},\uv{1}{2},\uv{3}{10},\uv{4}{9},\uv{5}{7},
\uv{6}{8},\uv{11}{16},\uv{13}{15},\uv{14}{17}\}.
\end{align*}
\endgroup
Each displayed set is a perfect matching, and
\[
  \bigcup_{j=1}^{4}M^{(a)}_j=E(\Ba),
  \qquad
  \bigcup_{j=1}^{4}M^{(b)}_j=E(\Bb).
\]

\enlargethispage{3\baselineskip}
\begin{thebibliography}{9}
\footnotesize
\setlength{\itemsep}{0.1em}

\bibitem{Blanusa1946}
D.~Blanu\v{s}a,
\emph{Problem \v{c}etiriju boja},
Glasnik Mat.-Fiz. Astr. Ser. II \textbf{1} (1946), 31--42.

\bibitem{FouquetVanherpe2009}
J.-L.~Fouquet and J.-M.~Vanherpe,
\emph{On the perfect matching index of bridgeless cubic graphs},
arXiv:0904.1296v1 (2009).

\bibitem{HagglundSteffen2014}
J.~H\"agglund and E.~Steffen,
\emph{Petersen-colorings and some families of snarks},
Ars Math. Contemp. \textbf{7} (2014), 161--173,
\doi{10.26493/1855-3974.288.11a}.

\bibitem{HakobyanMkrtchyan2019}
A.~Hakobyan and V.~Mkrtchyan,
\emph{$S_{12}$ and $P_{12}$-colorings of cubic graphs},
Ars Math. Contemp. \textbf{17} (2019), 431--445,
\doi{10.26493/1855-3974.1758.410}.

\bibitem{Putman2026}
B.~Putman,
\emph{A 112-vertex counterexample to the Petersen Coloring Conjecture},
arXiv:2608.10012v1, 8 August 2026.

\end{thebibliography}

\end{document}